\noindent\textbf{Can a visual language model preserve a change of meaning when it writes freely? We tested a bilingual model that reads printed glyph images and emits an answer image from a learned print repertoire. For 77 pairs of corpus-derived sentences drawn from its training pool, the writer scored the aligned translation above its competitor on 260 of 308 requests. Asked to write from the source alone, it conveyed the changed meaning under both source variants and in both languages for three of 77 pairs. The corresponding count was 46 of 51 for constructed event pairs. Across four unseen books the same writer earned six complete-meaning credits in 96 requests. These tests locate a gap between recognising which sentence fits a picture and composing the sentence without supplied options.}\par
\begin{multicols}{2}
\section*{Keeping the source edit visible}
Each pair shares an English sentence frame and differs by one word. We compared both corresponding Chinese sentences in both translation directions. The conditional score includes the learned stop. The free writer receives only the source image and derives its phrase draft from that image. We archived all 512 free answers. All stopped and all 32 sampled CPU replays reproduced the original pixels.
\end{multicols}
\displayfigure{Paired_visual_answers.pdf}{A visual writer can preserve a name and reverse the sense of a dream}{Original sources and free outputs for the first successful and first unsuccessful Chinese-to-English pair in a fixed order among corpus-derived pairs whose supplied translations passed independent model assessment. The writer changes Mom to Harry when the addressee changes. It writes \emph{nice dream} for a nightmare. The sampled sentences were in the broad training pool.}
\begin{multicols}{2}
\section*{Assessing the complete change}
For each source a quantized Qwen3-8B assessor saw both written answers in either order and could accept either, both or neither. A complete pair required acceptance of the aligned answer alone for both sources, in both orders, in both languages. Among corpus-derived requests, 125/308 individual answers received meaning credit in both orders and 21 matched the reference pixels exactly. Among constructed requests these counts were 200/204 and 103. No pair reproduced all four reference images exactly.
Of the 125 credited corpus-derived answers, 26 reproduced the complete phrase draft and 99 changed its pixels. Editing made 14 answers newly exact and spoiled 18 exact drafts.
\section*{Where a third glyph wins}
For the two candidate sentences we supplied their common prefix. At the first differing glyph, the writer preferred the aligned branch on 234/308 corpus-derived requests. Another glyph beat both branches on 112, including 83 where their relative ordering was already correct. Only 186 free outputs reached this canonical prefix. A preference between two complete answers can leave the actual next mark underdetermined.
\section*{Qualifying a training signal}
Of 1,384 mined pairs, two directional OPUS-MT models favoured the aligned answers by the required margin on 890. On the 128-family audit a separate assessor confirmed all four translations under both orders for 69, including only 21 of 77 corpus-derived pairs. The preregistered minimum was 112. We withheld these bulk labels. A byte-level trace identified the constructed examples in the original training mixture. \repo{references/paired_translation_teacher_result.md}{Provenance, scores, original answers and formal analysis}.
\end{multicols}
